\begin{abstract}
Appearance features are indispensable for online multi-object tracking (MOT), yet they are also the first cue to fail under heavy occlusion, crowding, and detector noise. Modern trackers therefore face a recurring question at association time: \emph{when should appearance be trusted, and when should the tracker fall back to motion?} We address this question with \textbf{Laplace-Guided Temporal Reliability Association (LTRA)}, a light-weight plug-in module built on top of a standard tracking-by-detection pipeline.

LTRA summarizes each track with \emph{multi-scale Laplace temporal signatures}, obtained by exponentially decayed aggregation of historical ReID features. It then estimates pairwise reliability from three factors: agreement between instantaneous appearance and temporal signatures, trajectory stability over time, and motion consistency. The resulting reliability score controls how strongly appearance cues influence the final association cost. LTRA does not modify the detector, does not require retraining the ReID encoder, and can be inserted into an existing strong tracker with minimal changes.

Our current proxy results already show that the proposed reliability mechanism improves association quality on MOT17 validation sequences: compared with the same baseline tracker, LTRA improves HOTA from 65.53 to 67.47, AssA from 57.71 to 60.60, and reduces ID switches from 223 to 167. The full paper is organized around a single claim: a simple and interpretable temporal reliability estimator can make appearance-based association substantially more robust without redesigning the entire tracking system.
\end{abstract}
